% ============================================================================
%  DogApp -- Miniproyecto 1 (Sprint 1)
%  Documento de especificacion tecnica de construccion derivado de las HU
%  y de la rubrica de evaluacion del curso
%  Curso: Desarrollo de Aplicaciones para Dispositivos Moviles
%  Compilar con: latexmk -pdf  (o pdflatex dos veces)
% ============================================================================
\documentclass[11pt,letterpaper]{article}

\usepackage[T1]{fontenc}
\usepackage[utf8]{inputenc}
\usepackage{lmodern}
\usepackage[spanish,es-tabla,es-noshorthands]{babel}

\usepackage[letterpaper,margin=2.3cm,top=2.6cm,bottom=2.5cm,headsep=0.55cm,footskip=1.0cm]{geometry}
\usepackage{microtype}
\usepackage{needspace}

\usepackage[table,dvipsnames]{xcolor}
\usepackage{graphicx}
\usepackage{booktabs}
\usepackage{xltabular}
\usepackage{tabularx}
\usepackage{array}
\usepackage{colortbl}
\usepackage{enumitem}
\usepackage{titlesec}
\usepackage{float}
\usepackage{tikz}
\usetikzlibrary{positioning,shapes.misc,calc,arrows.meta}
\usepackage[most]{tcolorbox}
\usepackage{listings}
\usepackage{fancyhdr}
\usepackage{pifont}
\usepackage{amssymb}
\usepackage[hidelinks]{hyperref}

% ---------------------------- Paleta DogApp --------------------------------
\definecolor{pink}{HTML}{F05861}
\definecolor{pinkdk}{HTML}{C9353F}
\definecolor{dark}{HTML}{2E2E33}
\definecolor{ink}{HTML}{1F2226}
\definecolor{muted}{HTML}{6B7280}
\definecolor{violet}{HTML}{6D3FCB}
\definecolor{okgreen}{HTML}{15803D}
\definecolor{badred}{HTML}{B91C1C}
\definecolor{amberx}{HTML}{B45309}
\definecolor{soft}{HTML}{FDF2F3}
\definecolor{softgray}{HTML}{F5F6F8}
\definecolor{linegray}{HTML}{D9DDE4}

\hypersetup{
  pdftitle={DogApp Sprint 1 - Especificacion de construccion},
  pdfauthor={Equipo DogApp - Desarrollo de Aplicaciones para Dispositivos Moviles},
  pdfsubject={Especificacion tecnica derivada de las Historias de Usuario y la rubrica},
  pdfkeywords={DogApp, MVVM, Room, Retrofit, Jira, Android, Sprint 1}
}

% ---------------------------- Estilo ---------------------------------------
\titleformat{\section}
  {\sffamily\Large\bfseries\color{pinkdk}}{\thesection}{0.7em}{}
  [\vspace{-0.5em}\color{pink}\rule{\linewidth}{1pt}]
\titleformat{\subsection}
  {\sffamily\large\bfseries\color{ink}}{\thesubsection}{0.6em}{}
\titleformat{\subsubsection}
  {\sffamily\normalsize\bfseries\color{violet}}{\thesubsubsection}{0.5em}{}
\titlespacing{\section}{0pt}{1.5em}{0.7em}
\titlespacing{\subsection}{0pt}{1.0em}{0.4em}

\pagestyle{fancy}
\fancyhf{}
\renewcommand{\headrulewidth}{0.4pt}
\renewcommand{\headrule}{\hbox to\headwidth{\color{linegray}\leaders\hrule height \headrulewidth\hfill}}
\fancyhead[L]{\footnotesize\sffamily\color{muted}DogApp $\cdot$ Sprint 1}
\fancyhead[R]{\footnotesize\sffamily\color{muted}Especificaci\'on de construcci\'on}
\fancyfoot[C]{\footnotesize\sffamily\color{muted}\thepage}

\setlength{\parskip}{0.45em}
\setlength{\parindent}{0pt}
\setlength{\emergencystretch}{3em}

\newcommand{\hex}[1]{\texttt{\##1}}
\newcommand{\code}[1]{\texttt{\small #1}}
\newcommand{\ok}{\textcolor{okgreen}{\ding{51}}}
\newcommand{\no}{\textcolor{badred}{\ding{55}}}
\newcommand{\warn}{\textcolor{amberx}{\ding{43}}}

% cajas reutilizables
\newtcolorbox{notabox}[1][]{
  colback=soft, colframe=pink, boxrule=0.6pt, arc=3pt,
  left=8pt, right=8pt, top=6pt, bottom=6pt,
  fonttitle=\sffamily\bfseries\small, title=#1
}
\newtcolorbox{techbox}[1][]{
  colback=softgray, colframe=linegray, boxrule=0.5pt, arc=3pt,
  left=8pt, right=8pt, top=6pt, bottom=6pt,
  fonttitle=\sffamily\bfseries\small\color{ink}, title=#1
}

\lstset{
  basicstyle=\ttfamily\footnotesize,
  breaklines=true,
  frame=l,
  framesep=6pt,
  rulecolor=\color{linegray},
  backgroundcolor=\color{softgray},
  xleftmargin=10pt,
  showstringspaces=false,
  columns=fullflexible,
  keepspaces=true,
  commentstyle=\color{muted}\itshape,
  keywordstyle=\color{violet}\bfseries,
  stringstyle=\color{pinkdk},
}

\renewcommand{\arraystretch}{1.25}
\renewcommand{\tabularxcolumn}[1]{>{\raggedright\arraybackslash}p{#1}}

\begin{document}

% ================================ PORTADA ==================================
\begin{titlepage}
\centering
\vspace*{0.6cm}

\begin{tikzpicture}[scale=1]
  \fill[dark, rounded corners=10pt] (0,0) rectangle (6.2,1.9);
  \node[white, font=\sffamily\bfseries\Large] at (3.1,1.05) {DogApp};
  \node[white, font=\sffamily\footnotesize] at (3.1,0.45) {Agenda de citas para mascotas};
\end{tikzpicture}

\vspace{1.4cm}
{\sffamily\color{muted}\large Universidad del Valle}\\[2pt]
{\sffamily\color{muted}Escuela de Ingenier\'ia de Sistemas y Computaci\'on $\cdot$ Sede Cali}\\[1.4cm]

{\sffamily\bfseries\Huge\color{dark} Documento de construcci\'on}\\[6pt]
{\sffamily\bfseries\Large\color{pinkdk} Miniproyecto 1 --- Sprint 1}\\[1.2cm]

{\large\itshape Qu\'e hay que construir, c\'omo, y c\'omo se gana cada d\'ecima\\
de la r\'ubrica de evaluaci\'on}

\vfill

\begin{tabularx}{0.86\linewidth}{@{}>{\sffamily\bfseries\color{muted}}l X@{}}
\toprule
Curso & Desarrollo de Aplicaciones para Dispositivos M\'oviles\\
Docente & Ing. Walter Medina\\
Proyecto & DogApp --- Agenda de citas para mascotas (5 Historias de Usuario)\\
Alcance & Sprint 1: HU 1.0 a HU 5.0 + r\'ubrica RA-1, RA-2, RA-3\\
Sprint & 11 de noviembre $\to$ 25 de noviembre de 2026, 6:30 p.m.\\
Monitor\'ia & 18 de noviembre de 2026\\
Entregable & PDF \code{equipoUno.\allowbreak{}pdf} (o el nombre del equipo) + repositorio + Jira\\
\bottomrule
\end{tabularx}

\vspace{1.2cm}
{\footnotesize\sffamily\color{muted}
Documento de trabajo del equipo. Elaborado a partir del enunciado oficial de las\\
Historias de Usuario y de la tabla de evaluaci\'on del Miniproyecto 1.}
\end{titlepage}

\tableofcontents
\newpage

% ============================ 1. RESUMEN ===================================
\section{Qu\'e se construye, en una p\'agina}

DogApp es una aplicaci\'on Android nativa para que el administrador de una
veterinaria registre y gestione las citas de las mascotas que debe atender.
El Sprint 1 entrega \textbf{cinco pantallas} y \textbf{siete capacidades
t\'ecnicas obligatorias}.

\subsection{Las cinco pantallas}

\begin{enumerate}[leftmargin=1.6em,itemsep=2pt]
  \item \textbf{Login (HU 1.0):} pantalla de bienvenida oscura con la marca
        DogApp y una huella animada; al tocarla se abre la autenticaci\'on
        biom\'etrica del dispositivo y, si el dedo es reconocido, se entra a la
        aplicaci\'on.
  \item \textbf{Home Administrador de Citas (HU 2.0):} lista de citas guardadas en
        una base de datos local, cada una con foto circular de la mascota,
        nombre, s\'intoma y turno; bot\'on flotante rosado para crear una cita.
  \item \textbf{Nueva Cita (HU 3.0):} formulario con cuatro campos validados y una
        lista de s\'intomas; el bot\'on Guardar solo se habilita cuando todo
        est\'a completo.
  \item \textbf{Detalle de la Cita (HU 4.0):} ficha del paciente con dos tarjetas
        superpuestas y dos botones flotantes: eliminar (rosado) y editar
        (azul oscuro).
  \item \textbf{Editar Cita (HU 5.0):} el mismo formulario, precargado, que
        actualiza el registro y regenera la foto seg\'un la nueva raza.
\end{enumerate}

\subsection{Las siete capacidades t\'ecnicas que exige la r\'ubrica}

La tabla de evaluaci\'on del enunciado lista exactamente estos siete items dentro
del RA-1 (60\,\%). Si falta uno, se pierde la fracci\'on correspondiente de ese
60\,\%:

\begin{center}
\begin{xltabular}{\linewidth}{@{}c X c@{}}
\toprule
\textbf{\#} & \textbf{Capacidad} & \textbf{D\'onde se evidencia}\\
\midrule
\endhead
\bottomrule
\endlastfoot
1 & Arquitectura MVVM \textbf{con Repository} & Paquetes \code{model}, \code{view}, \code{viewmodel}, \code{repository}, \code{data}\\
2 & Funcionamiento seg\'un las Historias de Usuario & Las 5 pantallas y sus criterios\\
3 & Base de datos SQLite con \textbf{Room} & Entidad, DAO, Database, Repository\\
4 & Consumo de API REST & Retrofit contra \code{dog.ceo}\\
5 & Uso de \textbf{Fragments} & 5 fragments, no 5 Activities\\
6 & Uso de \textbf{Navigation} & \code{nav\_graph.\allowbreak{}xml} con acciones y argumentos\\
7 & Interfaz id\'entica a los pantallazos (RA-2) & Sistema de dise\~no de la secci\'on~\ref{sec:diseno}\\
\end{xltabular}
\end{center}

\begin{notabox}[Fuera de alcance del Sprint 1 --- no lo construyas]
No hay autenticaci\'on con usuario y contrase\~na, ni registro, ni roles, ni
recuperaci\'on de clave. No hay notificaciones, ni calendario, ni recordatorios,
ni historial cl\'inico, ni sincronizaci\'on en la nube. No hay pantalla de
configuraci\'on. La autenticaci\'on es \emph{exclusivamente} biom\'etrica
dactilar, y la persistencia es \emph{exclusivamente} local con Room.
Invertir tiempo aqu\'i no suma nota y resta tiempo al sprint.
\end{notabox}

% ============================ 2. RUBRICA ===================================
\section{La r\'ubrica manda: c\'omo se reparte la nota}
\label{sec:rubrica}

La nota del equipo \textbf{no depende de que la app se vea bonita}: depende de
que cada item de la tabla se pueda demostrar en la sustentaci\'on. Esta es la
distribuci\'on textual del enunciado.

\begin{center}
\begin{xltabular}{\linewidth}{@{}>{\raggedright\arraybackslash}p{5.7cm} c >{\raggedright\arraybackslash}X@{}}
\toprule
\textbf{RA} & \textbf{\%} & \textbf{Items a calificar}\\
\midrule
\endhead
\bottomrule
\endlastfoot
\rowcolor{softgray}
RA-1. Implementa aplicaciones nativas para una plataforma m\'ovil & \textbf{60} &
Uso de arquitectura MVVM (con Repository) $\cdot$ Funcionamiento de la app como se
indica en las HU $\cdot$ Uso de base de datos SQLite con Room $\cdot$ Consumo de
API REST $\cdot$ Uso de Fragment $\cdot$ Uso de Navigation\\
RA-2. Usa principios de Dise\~no de Experiencia y de Interacci\'on & \textbf{20} &
Interfaz de la app como se indica en las Historias de Usuario\\
\rowcolor{softgray}
RA-3. Emplea gesti\'on de proyectos y control de versi\'on & \textbf{20} &
Gesti\'on de las tareas (criterios) de las HU en Jira: ``Por hacer, en curso,
hecho'' $\cdot$ Uso de GitHub donde se evidencie activamente el trabajo en
equipo durante el sprint\\
\midrule
\multicolumn{2}{@{}r}{\textbf{Nota equipo}} & \textbf{100\,\%}\\
\end{xltabular}
\end{center}

\subsection{La nota individual multiplica (no promedia)}

\begin{tcolorbox}[colback=soft,colframe=pink,boxrule=0.6pt,arc=3pt]
\centering\sffamily
\textbf{Nota definitiva del estudiante} $=$ \textbf{Nota equipo} $\times$ \textbf{Nota individual}\\[3pt]
{\footnotesize\color{muted}Ejemplo del enunciado: equipo 5 $\times$ individual 0,8 $=$ \textbf{4,0}}
\end{tcolorbox}

La nota individual va de 0 a 1 y el docente la obtiene de
\emph{``corroborar que cada uno de los integrantes particip\'o activamente''}.
En la pr\'actica eso se lee en tres lugares, y por eso el RA-3 no es tr\'amite:

\begin{itemize}[leftmargin=1.4em,itemsep=2pt]
  \item \textbf{Commits de GitHub con el nombre de cada integrante} (autor\'ia real,
        no un solo compa\~nero subiendo todo).
  \item \textbf{Movimiento de tarjetas en Jira} con responsable asignado.
  \item \textbf{Capacidad de cada uno de explicar y modificar su parte en vivo}
        durante la sustentaci\'on.
\end{itemize}

\noindent Un equipo con la app perfecta donde solo dos personas tienen commits
puede terminar con notas individuales de 0,5 para el resto. \textbf{El reparto de
tareas es parte del dise\~no t\'ecnico, no un adorno.}

\subsection{Penalizaciones expl\'icitas}

\begin{center}
\begin{tabularx}{\linewidth}{@{}X c@{}}
\toprule
\textbf{Situaci\'on} & \textbf{Efecto}\\
\midrule
No encender la c\'amara durante la sustentaci\'on & $-0{,}5$ en la nota final\\
Commit despu\'es de la fecha y hora de entrega & $-0{,}2$ en la nota final\\
Cualquier indicio de copia & Sanciones de los estatutos de la Universidad\\
Presentar de forma individual (la presentaci\'on es en equipo) & No admitido\\
\bottomrule
\end{tabularx}
\end{center}

\subsection{Mapa: de la r\'ubrica al artefacto concreto}

Esta es la tabla que conviene tener abierta mientras se trabaja. Cada fila es un
item calificable y el artefacto exacto que lo demuestra.

\begin{center}
\small
\begin{xltabular}{\linewidth}{@{}l l X@{}}
\toprule
\textbf{RA} & \textbf{Item} & \textbf{Artefacto que lo demuestra}\\
\midrule
\endhead
\bottomrule
\endlastfoot
RA-1 & MVVM + Repository & Estructura de paquetes + \code{CitaViewModel} observando \code{CitaRepository}; ninguna consulta a Room ni a Retrofit dentro de un Fragment\\
RA-1 & Funcionamiento HU & Recorrido completo: login $\to$ home $\to$ nueva $\to$ home $\to$ detalle $\to$ editar $\to$ home $\to$ detalle $\to$ borrar\\
RA-1 & Room & \code{CitaEntity}, \code{CitaDao}, \code{DogAppDatabase} y base persistente tras cerrar la app\\
RA-1 & API REST & \code{DogApiService} con Retrofit: lista de razas y foto por raza\\
RA-1 & Fragment & \code{LoginFragment}, \code{HomeFragment}, \code{NuevaCitaFragment}, \code{DetalleFragment}, \code{EditarFragment}\\
RA-1 & Navigation & \code{res/navigation/\allowbreak{}nav\_graph.\allowbreak{}xml} con 5 destinos, 6--8 acciones y argumento \code{citaId}\\
RA-2 & Interfaz & Comparaci\'on lado a lado con los 7 pantallazos del enunciado\\
RA-3 & Jira & Tablero con columnas Por hacer / En curso / Hecho e historias movi\'endose durante 2 semanas\\
RA-3 & GitHub & Repositorio \code{equipoUno}, ramas por HU, Pull Requests y commits de todos\\
\end{xltabular}
\end{center}

% ============================ 3. ARQUITECTURA ==============================
\section{Arquitectura obligatoria: MVVM + Repository}
\label{sec:mvvm}

Las diapositivas de la clase 9 fijan la definici\'on que el docente va a usar como
vara: \emph{Model} (datos, sin conocer la interfaz), \emph{View} (Activity o
Fragment, solo l\'ogica de interfaz), \emph{ViewModel} (conecta Model con View por
medio de LiveData y maneja la l\'ogica) y \emph{Repository} (puente entre las
fuentes de datos y el ViewModel). Sobre eso, la r\'ubrica agrega dos capacidades
observables: Room para SQLite y Retrofit para la API.

\subsection{El flujo que hay que poder dibujar en el tablero}

\begin{center}
\begin{tikzpicture}[
  node distance=0.75cm,
  box/.style={draw=linegray,fill=white,rounded corners=4pt,minimum height=0.85cm,
              align=center,font=\sffamily\footnotesize,inner sep=4pt,thick},
  core/.style={box,fill=soft,draw=pink},
  arr/.style={-{Stealth[length=2mm]},line width=0.7pt,draw=ink}
]
\node[core,text width=6.2cm] (view)
  {\textbf{VIEW} --- Fragments + XML\\Login $\cdot$ Home $\cdot$ Nueva $\cdot$ Detalle $\cdot$ Editar};
\node[box,text width=6.2cm,below=of view] (vm)
  {\textbf{VIEWMODEL} --- \code{LiveData}, \code{viewModelScope}\\l\'ogica de validaci\'on y de pantalla};
\node[core,text width=6.2cm,below=of vm] (repo)
  {\textbf{REPOSITORY} --- \code{CitaRepository}\\\'unica puerta de entrada a los datos};
\node[box,text width=2.85cm,below=1.0cm of repo.south west,anchor=north west] (room)
  {\textbf{ROOM}\\SQLite local\\\code{CitaDao}};
\node[box,text width=2.85cm,below=1.0cm of repo.south east,anchor=north east] (api)
  {\textbf{RETROFIT}\\API dog.ceo\\razas + fotos};

\draw[arr] (view) -- node[right,font=\sffamily\scriptsize,color=muted]{observa LiveData} (vm);
\draw[arr] ([xshift=-1.2cm]vm.south) -- node[left,font=\sffamily\scriptsize,color=muted]{pide datos} ([xshift=-1.2cm]repo.north);
\draw[arr] ([xshift=1.2cm]repo.north) -- node[right,font=\sffamily\scriptsize,color=muted]{entrega} ([xshift=1.2cm]vm.south);
\draw[arr] (repo.south) ++(-1.4,0) -- (room.north);
\draw[arr] (repo.south) ++(1.4,0) -- (api.north);
\end{tikzpicture}
\end{center}

\subsection{Reglas de oro que el jurado verifica}

\begin{enumerate}[leftmargin=1.6em,itemsep=3pt]
  \item \textbf{Ning\'un Fragment toca Room ni Retrofit.} Ni un \code{findViewById}
        de base de datos, ni un \code{Call} de Retrofit dentro de la vista. Todo
        pasa por el ViewModel y \'este por el Repository.
  \item \textbf{Ning\'un ViewModel conoce \code{Context}, \code{View} ni
        \code{Activity}.} Si aparece \code{Toast} dentro de un ViewModel, ya hay
        una fuga de capa.
  \item \textbf{Toda la UI se actualiza observando \code{LiveData}} (o
        \code{StateFlow} si el equipo domina corrutinas; el enunciado de la clase
        habla de LiveData, as\'i que LiveData es la apuesta segura).
  \item \textbf{Las operaciones de base de datos son \code{suspend}} y se invocan
        con \code{viewModelScope.\allowbreak{}launch}.
  \item \textbf{El Repository es la \'unica fuente de verdad.}
  \item \textbf{Una sola Activity} (\code{MainActivity}) con un
        \code{NavHostFragment}; las cinco pantallas son Fragments.
\end{enumerate}

\begin{notabox}[Se\~nal de alarma para la nota de arquitectura]
Si el profesor pide ``mu\'estrame el ViewModel de Nueva Cita'' y lo que aparece es
un Fragment de 300 l\'ineas con \code{db.citaDao(\allowbreak{}).insert(\allowbreak{}...)}, el item MVVM
vale cero aunque la app funcione. La app tiene que funcionar \emph{y} estar
partida en capas.
\end{notabox}

\newpage

\subsection{Paquetes y clases que se deben crear}

\begin{verbatim}
com.univalle.dogapp
|-- data/
|   |-- local/
|   |   |-- CitaEntity.kt          @Entity(tableName = "citas")
|   |   |-- CitaDao.kt             @Dao
|   |   `-- DogAppDatabase.kt      @Database(entities=[CitaEntity::class], version=1)
|   `-- remote/
|       |-- DogApiService.kt       @GET breeds/list/all, breed/{raza}/images/random
|       |-- BreedsResponse.kt      data class (message: Map<String, List<String>>)
|       `-- ImageResponse.kt       data class (message: String, status: String)
|-- repository/
|   `-- CitaRepository.kt          une DAO + API + Glide
|-- model/
|   `-- Cita.kt                    modelo de dominio / UI
|-- viewmodel/
|   |-- LoginViewModel.kt          estado de autenticacion
|   |-- HomeViewModel.kt           LiveData<List<Cita>>
|   |-- FormCitaViewModel.kt       crea y edita + razas + foto
|   `-- DetalleViewModel.kt        LiveData<Cita?> + eliminar
|-- view/
|   |-- MainActivity.kt            host del NavHostFragment
|   |-- login/LoginFragment.kt
|   |-- home/HomeFragment.kt  +  CitaAdapter.kt  +  CitaViewHolder.kt
|   |-- nueva/NuevaCitaFragment.kt
|   |-- detalle/DetalleFragment.kt
|   `-- editar/EditarFragment.kt
`-- utils/
    |-- Constants.kt               BASE_URL, nombres de argumentos
    `-- BiometricHelper.kt         envoltorio de BiometricPrompt
\end{verbatim}

\begin{techbox}[Atajo leg\'itimo: el proyecto de la clase 8]
El proyecto \code{clase8\_univalle} que acompa\~na el curso tiene \emph{exactamente}
la misma arquitectura que pide este miniproyecto, pero con un inventario:
\code{InventoryDao}, \code{InventoryDB}, \code{InventoryRepository},
\code{InventoryViewModel}, cuatro Fragments (Home, Add, Details, Edit),
\code{nav\_graph.\allowbreak{}xml}, RecyclerView con adapter y viewholder, y Room 2.5.2 con KSP.
Sirve como \textbf{esqueleto de referencia del docente}: renombrar a
\code{Cita}, cambiar el modelo de datos, a\~nadir Retrofit y la biometr\'ia, y
rehacer por completo la interfaz. Usar material del curso no es copia; copiar
el repositorio de otro equipo s\'i lo es.
\end{techbox}

% ============================ 4. ROOM ======================================
\section{Modelo de datos y persistencia con Room}
\label{sec:room}

\subsection{La entidad}

El enunciado pide cinco datos m\'as el turno. El turno \textbf{es el ID
autogenerado} de la base de datos, y la foto depende de la raza, as\'i que hay
que guardar tambi\'en la URL de la imagen.

\begin{center}
\small
\begin{xltabular}{\linewidth}{@{}l l X l@{}}
\toprule
\textbf{Campo} & \textbf{Tipo} & \textbf{Regla del enunciado} & \textbf{Nulo}\\
\midrule
\endhead
\bottomrule
\endlastfoot
\code{id} & \code{Int} & \code{@PrimaryKey(\allowbreak{}autoGenerate =\allowbreak{} true)}. Es el \textbf{turno} que se pinta como \code{\#1}, \code{\#2}\ldots\ en el Home y en el Detalle & \no\\
\code{nombreMascota} & \code{String} & Obligatorio, m\'aximo 15 caracteres & \no\\
\code{raza} & \code{String} & Obligatorio; se elige del autocompletado de razas & \no\\
\code{propietario} & \code{String} & Obligatorio, m\'aximo 30 caracteres & \no\\
\code{telefono} & \code{String} & Obligatorio, solo d\'igitos, m\'aximo 10 & \no\\
\code{sintoma} & \code{String} & Uno de los seis s\'intomas de la lista & \no\\
\code{imagenUrl} & \code{String?} & URL que devuelve la API para esa raza & s\'i\\
\code{fechaCreacion} & \code{Long} & \code{System.\allowbreak{}currentTimeMillis(\allowbreak{})}; desempata el orden & \no\\
\end{xltabular}
\end{center}

\begin{notabox}[Dos decisiones que evitan bugs cl\'asicos]
\textbf{(a)} Guarda \code{telefono} como \code{String}, no como \code{Int}:
si se guarda como entero, un n\'umero que empiece por 0 pierde el cero y un
tel\'efono con espacios revienta el parseo. \textbf{(b)} Guarda la
\code{raza} tal como se muestra (``Pug'', ``Pitbull''), pero para llamar a la API
normaliza en el Repository (\code{raza.trim().lowercase(\allowbreak{})}) porque el endpoint
espera la clave en min\'usculas. No guardes la clave de la API como texto visible
del usuario.
\end{notabox}

\subsection{DAO: las cinco operaciones del sprint}

\begin{verbatim}
@Dao
interface CitaDao {
    @Insert(onConflict = OnConflictStrategy.REPLACE)
    suspend fun insertar(cita: CitaEntity): Long      // devuelve el turno nuevo

    @Update
    suspend fun actualizar(cita: CitaEntity)

    @Delete
    suspend fun eliminar(cita: CitaEntity)

    @Query("SELECT * FROM citas ORDER BY id ASC")
    fun observarTodas(): LiveData<List<CitaEntity>>   // reactivo para el Home

    @Query("SELECT * FROM citas WHERE id = :id")
    fun observarPorId(id: Int): LiveData<CitaEntity?> // para Detalle y Editar
}
\end{verbatim}

\begin{notabox}[El detalle que rompe la mitad de los proyectos]
Si \code{observarTodas(\allowbreak{})} devuelve \code{List<CitaEntity>} en lugar de
\code{LiveData<.\allowbreak{}..>}, la lista del Home \textbf{no se refresca} despu\'es de
guardar o eliminar, y hay que llamar a mano al m\'etodo cada vez. Con
\code{LiveData}, la pantalla se actualiza sola y el criterio
``debe visualizarse la cita guardada'' (HU 3.0 y HU 5.0) se cumple sin c\'odigo
extra. Es tambi\'en la evidencia de ``LiveData'' de la clase 9.
\end{notabox}

\subsection{Base de datos y Repository}

\begin{verbatim}
@Database(entities = [CitaEntity::class], version = 1, exportSchema = false)
abstract class DogAppDatabase : RoomDatabase() {
    abstract fun citaDao(): CitaDao
    companion object {
        @Volatile private var INSTANCIA: DogAppDatabase? = null
        fun obtener(context: Context): DogAppDatabase =
            INSTANCIA ?: synchronized(this) {
                INSTANCIA ?: Room.databaseBuilder(
                    context.applicationContext,
                    DogAppDatabase::class.java, "dogapp.db"
                ).build().also { INSTANCIA = it }
            }
    }
}
\end{verbatim}

El \code{CitaRepository} recibe el DAO y el servicio de Retrofit y expone
funciones de negocio, no de base de datos:

\begin{verbatim}
class CitaRepository(private val dao: CitaDao, private val api: DogApiService) {

    fun citas(): LiveData<List<CitaEntity>> = dao.observarTodas()
    fun cita(id: Int): LiveData<CitaEntity?> = dao.observarPorId(id)

    suspend fun guardar(cita: CitaEntity): Long = dao.insertar(cita)
    suspend fun modificar(cita: CitaEntity) = dao.actualizar(cita)
    suspend fun borrar(cita: CitaEntity) = dao.eliminar(cita)

    suspend fun razas(): List<String> = withContext(Dispatchers.IO) {
        api.listarRazas().message.keys
            .map { it.replaceFirstChar { c -> c.uppercase() } }
            .sorted()
    }

    suspend fun fotoDeRaza(raza: String): String? = withContext(Dispatchers.IO) {
        runCatching { api.fotoPorRaza(raza.trim().lowercase()).message }.getOrNull()
    }
}
\end{verbatim}

\begin{techbox}[Inyecci\'on de dependencias sin frameworks]
No hace falta Hilt ni Koin. Basta una clase \code{ServiceLocator} (o un
\code{object AppContainer}) que construya la base de datos, el servicio Retrofit y
el repository una sola vez, y f\'abricas de ViewModel
(\code{ViewModelProvider.\allowbreak{}Factory}) que los reciban. La r\'ubrica no pide DI; s\'i
pide que las capas est\'en separadas.
\end{techbox}

% ============================ 5. API REST ==================================
\section{Consumo de API REST con Retrofit}
\label{sec:api}

La app consume dos servicios p\'ublicos de \code{dog.ceo}; ambos son GET y no
requieren llave.

\begin{center}
\small
\begin{xltabular}{\linewidth}{@{}l l X@{}}
\toprule
\textbf{Uso} & \textbf{Endpoint} & \textbf{Respuesta}\\
\midrule
\endhead
\bottomrule
\endlastfoot
Lista de razas para el autocompletado &
\url{https://dog.ceo/api/breeds/list/all} &
\code{message} es un objeto del tipo \code{\{"raza":\allowbreak{} ["subrazas"]\}}:
\code{pug} viene con lista vac\'ia y \code{bulldog} con
\code{["boston",\allowbreak{}"english"]}\\
Foto de una raza concreta &
\url{https://dog.ceo/api/breed/}\code{\{raza\}}\url{/images/random} &
\code{message} trae la URL de una foto de esa raza y \code{status} vale
\code{"success"}\\
\end{xltabular}
\end{center}

\begin{notabox}[\warn\ El enunciado sugiere \code{breeds/\allowbreak{}image/random}: le\'elo con cuidado]
El texto de la HU 3.0 y HU 5.0 dice que para la foto se debe consumir
\code{https://dog.ceo/api/breeds/\allowbreak{}image/random}, pero ese endpoint devuelve una
\textbf{foto de una raza aleatoria}, no de la raza elegida. Y el mismo enunciado
exige que ``si la raza escogida fue Akita, la imagen debe mostrar un perro alusivo
a esta raza''. Las dos frases no se cumplen con el mismo endpoint.
\textbf{Soluci\'on defendible:} usar \code{/breed/\{raza\}/\allowbreak{}images/\allowbreak{}random} para
cumplir el requisito funcional, y dejar \code{breeds/\allowbreak{}image/random} como
\emph{fallback} cuando la raza no exista en la API (por ejemplo si el usuario
escribe una raza a mano) o cuando falle la red. Menci\'onalo en el README y en la
sustentaci\'on: demuestra criterio, no error.
\end{notabox}

\subsection{Configuraci\'on de Retrofit}

\begin{verbatim}
interface DogApiService {
    @GET("api/breeds/list/all")
    suspend fun listarRazas(): BreedsResponse

    @GET("api/breed/{raza}/images/random")
    suspend fun fotoPorRaza(@Path("raza") raza: String): ImageResponse

    @GET("api/breeds/image/random")
    suspend fun fotoAleatoria(): ImageResponse
}

object RetrofitClient {
    private const val BASE_URL = "https://dog.ceo/"
    val api: DogApiService by lazy {
        val log = HttpLoggingInterceptor().apply { level = BODY }
        Retrofit.Builder()
            .baseUrl(BASE_URL)
            .addConverterFactory(GsonConverterFactory.create())
            .client(OkHttpClient.Builder().addInterceptor(log).build())
            .build()
            .create(DogApiService::class.java)
    }
}
\end{verbatim}

\subsection{Reglas de consumo que evitan perder puntos}

\begin{enumerate}[leftmargin=1.6em,itemsep=3pt]
  \item \textbf{Una sola llamada a la lista de razas por sesi\'on}, no una por
        tecla. Se carga en el \code{init} del ViewModel, se expone como
        \code{LiveData<List<String>>} y el adapter filtra en memoria.
  \item \textbf{El filtro del autocompletado es \code{startsWith}}, no
        \code{contains}: el criterio dice ``que comiencen por esas letras''.
        Umbral: 2 caracteres (\code{android:\allowbreak{}completionThreshold=\allowbreak{}"2"}).
  \item \textbf{Todo en corrutinas}: \code{suspend} + \code{Dispatchers.\allowbreak{}IO}. Nada
        de red en el hilo principal (NetworkOnMainThreadException en 2 segundos).
  \item \textbf{Tolera el fallo}: si no hay internet, la app tiene que seguir
        creando citas con una imagen de marcador. Envuelve con
        \code{runCatching} y no dejes que una excepci\'on cierre la pantalla.
  \item \textbf{Rotaci\'on}: un \code{Retrofit} por proceso, no uno por Fragment.
  \item \textbf{Manifiesto}: \code{<uses-permission android:\allowbreak{}name="android.\allowbreak{}permission.\allowbreak{}INTERNET"/\allowbreak{}>}.
  \item \textbf{Carga de im\'agenes con Glide} (o Picasso/Coil):
        \code{Glide.with(view).load(url)} con \code{placeholder(\allowbreak{}R.drawable.\allowbreak{}ic\_dog\_placeholder)} y \code{into(imageView)}
        evita bloqueos, memoriza y recorta al c\'irculo con
        \code{circleCrop(\allowbreak{})}.
\end{enumerate}

% ============================ 6. DISENO ====================================
\section{Sistema de dise\~no extra\'ido de los pantallazos}
\label{sec:diseno}

El enunciado advierte que no hay prototipo con medidas exactas y que hay que
parecerse ``lo m\'as posible'' a los pantallazos. Estos son los valores que se
leen de las capturas y que hay que centralizar en \code{res/values} para que las
cinco pantallas queden consistentes (eso es el RA-2).

\subsection{Tokens de color}

\begin{center}
\small
\begin{xltabular}{\linewidth}{@{}l l X@{}}
\toprule
\textbf{Token} & \textbf{Valor} & \textbf{Uso}\\
\midrule
\endhead
\bottomrule
\endlastfoot
\code{dogapp\_pink} & \hex{F05861} & Botones Guardar/Editar, FAB de crear, flechas de atr\'as, turno, l\'inea divisoria del detalle, FAB de eliminar\\
\code{dogapp\_dark} & \hex{2E2E33} & Fondo del login y de todos los headers/toolbars\\
\code{dogapp\_pink\_disabled} & \hex{F5A9AD} & Estado inactivo del bot\'on Guardar Cita\\
\code{dogapp\_field\_focus} & \hex{6D3FCB} & Borde del campo cuando tiene el foco (morado)\\
\code{dogapp\_edit\_blue} & \hex{2B2B6B} & FAB de editar en el detalle\\
\code{white} & \hex{FFFFFF} & Texto sobre oscuro, fondo de tarjetas, fondo de formularios\\
\code{hint\_gray} & \hex{9E9E9E} & Hints de los campos y textos secundarios\\
\code{ink} & \hex{212121} & Nombres, razas y contenido principal\\
\end{xltabular}
\end{center}

\subsection{Formas, elevaciones y tipograf\'ia}

\begin{center}
\small
\begin{xltabular}{\linewidth}{@{}l l l X@{}}
\toprule
\textbf{Elemento} & \textbf{Radio} & \textbf{Elevaci\'on} & \textbf{Nota}\\
\midrule
\endhead
\bottomrule
\endlastfoot
Header del Home & 48\,dp abajo-izquierda & 4\,dp & Solo la esquina inferior izquierda curvada; la derecha queda recta\\
Toolbar Nueva/Detalle/Editar & 28\,dp abajo (ambas) & 4\,dp & Esquinas inferiores curvadas, ver pantallazos\\
Tarjeta de \'item (Home) & 16\,dp & 4\,dp & Fondo blanco sobre fondo blanco: la separaci\'on la da la elevaci\'on y un margen de 8\,dp\\
Tarjeta de imagen (Detalle) & 16\,dp & 8\,dp & Va \emph{encima} de la tarjeta de datos\\
Tarjeta de datos (Detalle) & 16\,dp & 4\,dp & Se solapa por debajo de la de la imagen\\
Bot\'on Guardar/Editar & 24\,dp (c\'apsula) & 0--2\,dp & Rosa, texto blanco bold cuando est\'a habilitado\\
Campo de texto & 8\,dp & 0 & \code{OutlinedBox}, borde gris 1\,dp y morado 2\,dp con foco\\
FAB & 16\,dp & 6\,dp & 56\,dp; crear (abajo-der), borrar (abajo-izq), editar (abajo-der)\\
Avatar circular & 50\,\% & 0 & 52\,dp, \code{circleCrop}\\
\end{xltabular}
\end{center}

\subsection{Fuentes}

\begin{itemize}[leftmargin=1.4em,itemsep=2pt]
  \item \textbf{Marca ``DogApp'':} fuente externa redondeada y muy bold, como
        \emph{Baloo 2 ExtraBold}, \emph{Fredoka One} o \emph{Nunito ExtraBold}.
        Se descarga el \code{.ttf}, se copia en \code{res/font/dogapp\_brand.\allowbreak{}ttf}
        y se usa con \code{android:\allowbreak{}fontFamily=\allowbreak{}"@font/\allowbreak{}dogapp\_brand"}. Cumple el
        criterio ``no olvide usar fuentes externas que imiten ese estilo''.
  \item \textbf{T\'itulos de pantalla y contenido:} la sans del sistema
        (\code{sans-serif-medium} / \code{sans-serif-black}) es aceptable; si se
        quiere uniformidad, \emph{Poppins SemiBold} para t\'itulos.
  \item \textbf{Tama\~nos de referencia:} marca del login 40--44\,sp; subt\'itulo
        del Home 12\,sp; t\'itulo del header 20--22\,sp; nombre de mascota en el
        \'item 16\,sp bold; s\'intoma 13--14\,sp; turno \code{\#n} 12\,sp bold.
\end{itemize}

\subsection{Archivos de recursos a crear}

\begin{verbatim}
res/
|-- values/colors.xml            todos los tokens de arriba
|-- values/dimens.xml            spacing_8, radius_card_16, avatar_52 ...
|-- values/strings.xml           textos + array de sintomas
|-- values/themes.xml            Theme.DogApp (NoActionBar)
|-- font/dogapp_brand.ttf        fuente externa de la marca
|-- drawable/
|   |-- bg_header_home.xml       fondo oscuro, esquina inf-izq curva
|   |-- bg_toolbar_curved.xml    fondo oscuro, esquinas inf. curvadas
|   |-- bg_button_pink.xml       (o usar MaterialButton con cornerRadius)
|   |-- ic_back_pink.xml         flecha de atr\'as tintada de rosa
|   |-- ic_save.xml  ic_edit.xml  ic_delete.xml  ic_add.xml
|   `-- ic_dog_head.webp         cabeza de perro, fondo transparente
|-- color/til_stroke.xml         selector morado(focus) / gris(reposo)
|-- assets/fingerprint.json      animaci\'on Lottie de la huella
`-- mipmap-*/ic_launcher.*       \'icono propio dogApp (adaptive icon)
\end{verbatim}

% ============================ 7. PANTALLAS =================================
\section{Especificaci\'on pantalla por pantalla}
\label{sec:pantallas}

Cada subsecci\'on traduce los criterios del enunciado a implementaci\'on concreta.
La columna ``d\'onde'' es el archivo que el profesor pedir\'a abrir si algo no se
ve bien.

\subsection{HU 1.0 --- Ventana Login}

\textbf{Acta:} como aplicaci\'on, quiero presentar un sistema de logueo con
biometr\'ia dactilar para mejorar la experiencia de ingreso. \textbf{Tipo:}
Fragment de inicio, sin toolbar.

\begin{center}
\footnotesize
\begin{xltabular}{\linewidth}{@{}c X X@{}}
\toprule
\textbf{C} & \textbf{Qu\'e exige} & \textbf{C\'omo se implementa}\\
\midrule
\endhead
\bottomrule
\endlastfoot
1 & Fondo gris oscuro, sin toolbar & \code{Theme.DogApp} con \code{windowNoTitle=\allowbreak{}true} y \code{windowActionBar=\allowbreak{}false}; ra\'iz del layout con \code{android:\allowbreak{}background=\allowbreak{}"@color/\allowbreak{}dogapp\_dark"}. Confirmar que el fragment no infla \code{MaterialToolbar}\\
2 & Cabeza de perro arriba a la derecha, fondo transparente, solo la cabeza, optimizada & \code{ImageView} de un \code{.webp} con alfa recortado a la cabeza; anclado \code{topToTop} + \code{endToEnd} con m\'argen 24\,dp; \code{scaleType=\allowbreak{}"fitCenter"} y \code{layout\_constraintWidth\_percent=\allowbreak{}"0.55"} para que escale en pantallas distintas. Exportar en \code{drawable-xxhdpi} y \code{drawable-xxxhdpi} (o vectorizar)\\
3 & ``DogApp'' blanco bold, debajo de la imagen, centrado, fuente externa & \code{TextView} con \code{@font/dogapp\_brand}, \code{textColor=\allowbreak{}"@color/\allowbreak{}white"}, \code{textStyle=\allowbreak{}"bold"}, 40--44\,sp; \code{topToBottomOf} la imagen\\
4 & Imagen din\'amica de huella abajo al centro & \code{LottieAnimationView} con \code{lottie\_fileName=\allowbreak{}"fingerprint.\allowbreak{}json"} desde \code{assets/\allowbreak{}}, \code{autoPlay=\allowbreak{}"true"}, \code{repeatCount=\allowbreak{}"infinite"}; 120--160\,dp, anclada abajo con m\'argen inferior 64\,dp\\
5 & Al tocar la huella: ventana emergente con t\'itulo, subt\'itulo y ``Cancelar'' & \code{BiometricPrompt}. T\'itulo \code{"Autenticaci\'on con Biometr\'ia"}, subt\'itulo \code{"Ingrese su huella digital"}, \code{setNegativeButtonText(\allowbreak{}"Cancelar")}. El aspecto lo dibuja el sistema y \emph{var\'ia entre dispositivos}: el enunciado lo permite expl\'icitamente\\
6 & Huella correcta $\to$ HU 2.0; incorrecta $\to$ mensaje y no entra & En \code{onAuthenticationSucceeded}: \code{navigate(\allowbreak{}R.id.action\_login\_to\_home)} con \code{popUpTo=\allowbreak{}login} \code{inclusive=\allowbreak{}true} (para que ``atr\'as'' no vuelva al login). En \code{onAuthenticationError}: \code{Snackbar} con el mensaje del sistema y \emph{ninguna} navegaci\'on\\
7 & \'Icono propio de la app llamado dogApp & Adaptive icon: \code{mipmap-anydpi-v26/\allowbreak{}ic\_launcher.\allowbreak{}xml} con \code{foreground} (cabeza de perro) + \code{background} de color; \code{ic\_launcher\_round}; \code{android:\allowbreak{}label="dogApp"} en el manifest\\
\end{xltabular}
\end{center}

\begin{verbatim}
Manifest: <uses-permission android:name="android.permission.USE_BIOMETRIC"/>
Gradle:   implementation("androidx.biometric:biometric:1.1.0")

val prompt = BiometricPrompt(requireActivity(),
    ContextCompat.getMainExecutor(requireContext()),
    object : BiometricPrompt.AuthenticationCallback() {
        override fun onAuthenticationSucceeded(
            result: BiometricPrompt.AuthenticationResult) {
            findNavController().navigate(R.id.action_login_to_home)
        }
        override fun onAuthenticationError(code: Int, msg: CharSequence) {
            Snackbar.make(binding.root, msg, Snackbar.LENGTH_LONG).show()
        }
    })

val info = BiometricPrompt.PromptInfo.Builder()
    .setTitle("Autenticacion con Biometria")
    .setSubtitle("Ingrese su huella digital")
    .setNegativeButtonText("Cancelar")
    .setAllowedAuthenticators(BiometricManager.Authenticators.BIOMETRIC_STRONG)
    .build()
prompt.authenticate(info)
\end{verbatim}

\begin{notabox}[Prueba de la huella en el emulador]
Sin dedo enrolado, el di\'alogo no aparece y parece que ``la biometr\'ia no
funciona''. En el emulador: \emph{Settings $\to$ Security $\to$ Fingerprint} y
enrolar un dedo; luego, con la app en el di\'alogo, mandar
\code{adb -e emu finger touch 1} (o usar \emph{Extended Controls $\to$
Fingerprint}). Verifica tambi\'en el camino de error (huella no reconocida) y el
bot\'on Cancelar: los tres son criterios calificables.
\end{notabox}

\subsection{HU 2.0 --- Home Administrador de Citas}

\textbf{Acta:} como administrador, quiero visualizar las citas creadas para
gestionar su atenci\'on. \textbf{Tipo:} Fragment con RecyclerView + FAB.

\begin{center}
\footnotesize
\begin{xltabular}{\linewidth}{@{}c X X@{}}
\toprule
\textbf{C} & \textbf{Qu\'e exige} & \textbf{C\'omo se implementa}\\
\midrule
\endhead
\bottomrule
\endlastfoot
1 & El bot\'on atr\'as del tel\'efono debe llevar al escritorio, no al login &
\code{requireActivity(\allowbreak{}).onBackPressedDispatcher.\allowbreak{}addCallback(\allowbreak{}viewLifecycleOwner)\{ requireActivity(\allowbreak{}).finishAffinity(\allowbreak{}) \}}. Adem\'as, el login se elimina del back stack al navegar (ver HU 1.0, criterio 6)\\
2 & Fondo blanco; header gris oscuro con la esquina inferior izquierda curvada; t\'itulo blanco bold ``Administrador de Citas'' a la izquierda y centrado vertical; a la derecha la imagen del perro con ``DogApp'' debajo & \code{LinearLayout}/\code{ConstraintLayout} con \code{background=\allowbreak{}"@drawable/\allowbreak{}bg\_header\_home"} (shape s\'olido \hex{2E2E33} con \code{bottomLeftRadius=\allowbreak{}48dp} y \code{bottomRightRadius=\allowbreak{}0dp}). Texto en 2 l\'ineas a la izquierda; a la derecha \code{ImageView} de la cabeza + \code{TextView} ``DogApp'' en 12\,sp debajo\\
3 & \'Items: fondo blanco, bordes redondeados, elevaci\'on; imagen circular de la mascota a la izquierda; nombre y s\'intoma al centro; turno arriba a la derecha & \code{item\_cita.\allowbreak{}xml}: \code{MaterialCardView} (radio 16\,dp, elevaci\'on 4\,dp) con \code{ShapeableImageView} circular 52\,dp, \code{tvNombre} bold 16\,sp + \code{tvSintoma} 13\,sp gris, y \code{tvTurno} con texto \code{"\#"+id} en rosa bold arriba a la derecha\\
4 & FAB rosado \hex{F05861} abajo a la derecha que lleva a Nueva Cita & \code{FloatingActionButton} con \code{app:backgroundTint=\allowbreak{}"@color/\allowbreak{}dogapp\_pink"}, \code{srcCompat=\allowbreak{}"@drawable/\allowbreak{}ic\_add"}, \code{layout\_gravity=\allowbreak{}"bottom|end"}; \code{setOnClickListener \{ navigate(\allowbreak{}R.id.action\_home\_to\_nueva) \}}\\
5 & Al tocar un \'item, ir al Detalle de esa cita & \code{adapter.\allowbreak{}setOnItemClick \{ navigate(\allowbreak{}R.id.action\_home\_to\_detalle,\allowbreak{} bundleOf(\allowbreak{}"citaId" to it.\allowbreak{}id)) \}} o Safe Args\\
6 & La lista sale de Room (SQLite) & \code{viewModel.\allowbreak{}citas.observe(\allowbreak{}viewLifecycleOwner) \{ adapter.\allowbreak{}submitList(\allowbreak{}it) \}} --- \code{LiveData} generado por el DAO\\
\end{xltabular}
\end{center}

\begin{techbox}[C\'omo pintar la esquina curvada del header sin pelear con XML]
Dos caminos v\'alidos: \textbf{(a)} un \code{shape} drawable con
\code{<corners android:\allowbreak{}bottomLeftRadius=\allowbreak{}"48dp" android:\allowbreak{}bottomRightRadius=\allowbreak{}"0dp"/>}
y \code{<solid android:\allowbreak{}color="@color/\allowbreak{}dogapp\_dark"/\allowbreak{}>}; \textbf{(b)} un
\code{MaterialCardView}/\code{ImageView} con
\code{shapeAppearanceOverlay} usando \code{cornerFamily=\allowbreak{}rounded} y
\code{cornerSizeBottomLeft=\allowbreak{}48dp}, \code{cornerSizeBottomRight=\allowbreak{}0dp}. La opci\'on
(b) es la misma que se reutiliza en las toolbars de las HU 3.0, 4.0 y 5.0, as\'i
que conviene resolverla una vez y reutilizar el estilo.
\end{techbox}

\subsection{HU 3.0 --- Ventana Nueva Cita}

\textbf{Acta:} como administrador, quiero crear una cita para llevar el registro
de los pacientes. \textbf{Tipo:} Fragment de formulario, con toolbar personalizada.

\begin{center}
\footnotesize
\begin{xltabular}{\linewidth}{@{}c X X@{}}
\toprule
\textbf{C} & \textbf{Qu\'e exige} & \textbf{C\'omo se implementa}\\
\midrule
\endhead
\bottomrule
\endlastfoot
1 & Toolbar oscura con esquinas inferiores curvadas, t\'itulo ``Nueva Cita'' centrado blanco bold, flecha de atr\'as rosada que vuelve al Home & \code{MaterialToolbar} con \code{background=\allowbreak{}"@drawable/\allowbreak{}bg\_toolbar\_curved"}, \code{app:navigationIcon=\allowbreak{}"@drawable/\allowbreak{}ic\_back\_pink"} y \code{app:title} centrado (\code{app:titleCentered=\allowbreak{}"true"} o \code{TextView} centrado). \code{setNavigationOnClickListener \{ findNavController(\allowbreak{}).navigateUp(\allowbreak{}) \}}\\
2 & Campo ``Nombre de la Mascota'': hint gris, flota al enfocar, borde morado al enfocar, m\'aximo 15 caracteres & \code{TextInputLayout} estilo \code{Widget.\allowbreak{}MaterialComponents.\allowbreak{}TextInputLayout.\allowbreak{}OutlinedBox} + \code{TextInputEditText}, \code{hint="@string/\allowbreak{}nombre\_mascota"}, \code{android:\allowbreak{}maxLength=\allowbreak{}"15"}, \code{app:boxStrokeColor=\allowbreak{}"@color/\allowbreak{}til\_stroke"}. El flotado del hint es nativo del \code{TextInputLayout}\\
3 & Campo ``Raza'' con \textbf{autocompletado}: al escribir 2 letras aparecen las razas que \emph{comienzan} por ellas, alimentado por la API con Retrofit & \code{MaterialAutoCompleteTextView} dentro de \code{TextInputLayout} (\code{endIconMode=\allowbreak{}"dropdown\_menu"}), \code{android:\allowbreak{}completionThreshold=\allowbreak{}"2"}. El adapter se llena con \code{viewModel.\allowbreak{}razas.observe} y su \code{Filter} usa \code{startsWith(\allowbreak{}prefix,\allowbreak{} ignoreCase=\allowbreak{}true)}\\
4 & Campo ``Nombre del Propietario'', m\'aximo 30 & Igual que el criterio 2, \code{android:\allowbreak{}maxLength=\allowbreak{}"30"}\\
5 & Campo ``Tel\'efono'', solo n\'umeros, m\'aximo 10 & \code{android:\allowbreak{}inputType=\allowbreak{}"number"}, \code{android:\allowbreak{}maxLength=\allowbreak{}"10"}, \code{digits=\allowbreak{}"0123456789"}\\
6 & Lista de s\'intomas con el \'item ``S\'intomas'' por defecto (quemada en la app) & \code{MaterialAutoCompleteTextView} no editable (\code{inputType=\allowbreak{}"none"}, \code{endIconMode=\allowbreak{}"dropdown\_menu"}) con \code{ArrayAdapter} de \code{R.array.sintomas}; el \'indice 0 es ``S\'intomas'' = sin selecci\'on. Los seis valores: Solo duerme, No come, Fractura extremidad, Tiene pulgas, Tiene garrapatas, Bota demasiado pelo\\
7 & Bot\'on rosado, redondeado, con \'icono a la izquierda, texto ``Guardar Cita''; \textbf{inactivo} mientras falte nombre, raza, propietario o tel\'efono & \code{MaterialButton} con \code{app:cornerRadius=\allowbreak{}"24dp"}, \code{backgroundTint=\allowbreak{}"@color/\allowbreak{}dogapp\_pink"}, \code{app:icon="@drawable/\allowbreak{}ic\_save"}, \code{iconGravity=\allowbreak{}"textStart"}. Estado en \code{doAfterTextChanged}: \code{btn.isEnabled =\allowbreak{} nombre.\allowbreak{}isNotBlank(\allowbreak{}) \&\& raza.\allowbreak{}isNotBlank(\allowbreak{}) \&\& propietario.\allowbreak{}isNotBlank(\allowbreak{}) \&\& telefono.\allowbreak{}isNotBlank(\allowbreak{})}\\
8 & Si todo est\'a lleno pero no hay s\'intoma: al pulsar, mensaje ``Selecciona un s\'intoma'' & En el \code{onClick}: \code{if (sintomaSeleccionado =\allowbreak{}= null) \{ Snackbar/\allowbreak{}AlertDialog(\allowbreak{}"Selecciona un s\'intoma"); return \}}\\
9 & El bot\'on se habilita cuando los cuatro campos est\'an llenos y el texto pasa a blanco bold & Adem\'as de \code{isEnabled}, cambiar \code{setTextColor(\allowbreak{}Color.WHITE)} + \code{Typeface.\allowbreak{}BOLD} al habilitar, y \code{alpha=0.5} + color rosado claro mientras est\'e inactivo\\
10 & Al guardar: insertar en Room y navegar al Home, donde debe verse la cita & \code{viewModel.\allowbreak{}guardar(\allowbreak{}cita)} en \code{viewModelScope}; al terminar (\code{onSuccess}), \code{navigate(\allowbreak{}R.id.action\_nueva\_to\_home)} con \code{popUpTo} para no dejar el formulario en la pila\\
\end{xltabular}
\end{center}

\begin{notabox}[El filtro del autocompletado: el error m\'as com\'un del proyecto]
El \code{ArrayAdapter} por defecto filtra por ``contiene'' y respeta el orden de
la lista, as\'i que escribir ``pu'' puede devolver ``bullpug'' antes que ``pug''.
Para cumplir ``que comiencen por esas letras'' hay que sobrescribir el
\code{Filter} con la l\'ogica que sigue.
\end{notabox}

\begin{verbatim}
override fun getFilter(): Filter = object : Filter() {
    override fun performFiltering(prefix: CharSequence?): FilterResults {
        val p = prefix?.toString()?.trim()?.lowercase().orEmpty()
        val out = if (p.length < 2) emptyList()
                  else razas.filter { it.lowercase().startsWith(p) }
        return FilterResults().apply { values = out; count = out.size }
    }
    override fun publishResults(c: CharSequence?, r: FilterResults?) {
        clear()
        addAll(r?.values as? List<String> ?: emptyList())
        notifyDataSetChanged()
    }
}
\end{verbatim}

\noindent Y no olvides que \code{performFiltering} corre en hilo de fondo: no
llames a la API ah\'i. Pide la lista una sola vez y filtra en memoria.

\subsection{HU 4.0 --- Ventana Detalle de la Cita}

\textbf{Acta:} como administrador, quiero ver el detalle de un paciente para
decidir. \textbf{Tipo:} Fragment de solo lectura + dos FAB.

\begin{center}
\footnotesize
\begin{tabularx}{\linewidth}{@{}c X X@{}}
\toprule
\textbf{C} & \textbf{Qu\'e exige} & \textbf{C\'omo se implementa}\\
\midrule
1 & Toolbar oscura curvada; el t\'itulo es el \textbf{nombre de la mascota}; flecha rosada que vuelve al Home & Reutilizar el estilo de la HU 3.0; \code{toolbar.\allowbreak{}title =\allowbreak{} cita.nombreMascota} dentro del \code{observe}. Como el t\'itulo depende de datos, hay que asignarlo despu\'es de cargar el \code{LiveData}\\
2 & Card con la foto de la mascota, con margen, \textbf{por encima} de otra tarjeta blanca & \code{MaterialCardView} (radio 16\,dp, elevaci\'on 8\,dp) con \code{ShapeableImageView} (\code{scaleType=\allowbreak{}"centerCrop"}) y \code{layout\_margin=\allowbreak{}"20dp"}. Para superponer, \code{ConstraintLayout} con la segunda card anclada debajo y un \code{layout\_marginTop=\allowbreak{}"-32dp"} en la card de informaci\'on\\
3 & Segunda card: turno arriba a la derecha, raza en negro bold, s\'intoma, l\'inea delgada rosada, ``Propietario'' + nombre, ``Tel\'efono'' + n\'umero & Orden vertical: \code{tvTurno} (\code{"\#"+id}, rosa, bold, \code{end|top}), \code{tvRaza} (negro bold, centrado), \code{tvSintoma}, \code{View} de 1\,dp con \code{background=\allowbreak{}"@color/\allowbreak{}dogapp\_pink"} y ancho 90\,\%, \code{tvPropietario}, \code{tvTelefono}\\
4 & FAB rosado con \'icono de eliminaci\'on, abajo a la \textbf{izquierda}: borra en Room y vuelve al Home sin la cita & \code{FloatingActionButton} \code{layout\_gravity=\allowbreak{}"bottom|start"}, \code{backgroundTint=\allowbreak{}"@color/\allowbreak{}dogapp\_pink"}, \code{ic\_delete}. \code{viewModel.\allowbreak{}eliminar(\allowbreak{}cita); findNavController(\allowbreak{}).navigate(\allowbreak{}R.id.action\_detalle\_to\_home)}. No se exige confirmaci\'on\\
5 & FAB azul oscuro con \'icono de edici\'on, abajo a la \textbf{derecha}: lleva a Editar Cita & \code{FloatingActionButton} \code{layout\_gravity=\allowbreak{}"bottom|end"}, \code{backgroundTint=\allowbreak{}"@color/\allowbreak{}dogapp\_edit\_blue"}, \code{ic\_edit}; navega pasando el mismo \code{citaId}\\
\bottomrule
\end{tabularx}
\end{center}

\begin{verbatim}
// el detalle se alimenta del id que llega por navegacion
private val citaId: Int by lazy { requireArguments().getInt("citaId") }
viewModel.cargar(citaId)
viewModel.cita.observe(viewLifecycleOwner) { cita ->
    cita ?: run { findNavController().navigateUp(); return@observe }
    toolbar.title = cita.nombreMascota
    tvRaza.text = cita.raza
    tvSintoma.text = cita.sintoma
    tvTurno.text = "#${cita.id}"
    tvPropietario.text = "Propietario: ${cita.propietario}"
    tvTelefono.text = "Telefono: ${cita.telefono}"
    Glide.with(this).load(cita.imagenUrl).circleCrop().into(imgMascota)
}
\end{verbatim}

\begin{notabox}[La superposici\'on de tarjetas es el criterio m\'as f\'acil de perder]
El texto dice ``este contenedor est\'a por encima de otra tarjeta de color
blanco''. Dos \code{MaterialCardView} apiladas sin m\'as no lo logran: la de la
imagen debe tener mayor \code{elevation} \textbf{y} solaparse f\'isicamente con la
de abajo (margen superior negativo, o anclando la de arriba con
\code{layout\_marginBottom=\allowbreak{}"-32dp"}). Verif\'icalo en pantalla: si se ven dos
rect\'angulos separados con un hueco, el criterio no est\'a cumplido.
\end{notabox}

\subsection{HU 5.0 --- Ventana Editar Cita}

\textbf{Acta:} como administrador, quiero modificar datos de la mascota para
mantener actualizada la agenda. \textbf{Tipo:} mismo formulario de la HU 3.0,
precargado.

\begin{center}
\footnotesize
\begin{xltabular}{\linewidth}{@{}c X X@{}}
\toprule
\textbf{C} & \textbf{Qu\'e exige} & \textbf{C\'omo se implementa}\\
\midrule
\endhead
\bottomrule
\endlastfoot
1 & Toolbar ``Editar Cita'', flecha rosada que vuelve al \textbf{Detalle} (mostrando la mascota editada) & \code{navigateUp(\allowbreak{})} (no navegar al Home). El detalle se refresca solo porque observa \code{LiveData} de Room\\
2 & Nombre de la mascota, m\'aximo 15, \textbf{precargado} desde el detalle & \code{TextInputEditText.\allowbreak{}setText(\allowbreak{}cita.nombreMascota)} en el \code{observe}; no volver a pedir el id al usuario: llega por argumento de navegaci\'on\\
3 & Raza con el mismo autocompletado por API (2 letras, prefijo) & Reutilizar el adapter de la HU 3.0 (mismo c\'odigo, mismo ViewModel de formulario)\\
4 & Nombre del propietario, m\'aximo 30, precargado & \code{setText(\allowbreak{}cita.propietario)}\\
5 & Tel\'efono, solo n\'umeros, m\'aximo 10, precargado & \code{setText(\allowbreak{}cita.telefono)}\\
6 & Bot\'on rosado redondeado con \'icono y texto ``Editar Cita''; inactivo si falta alg\'un campo & Misma l\'ogica de habilitaci\'on que la HU 3.0, con \code{app:icon="@drawable/\allowbreak{}ic\_edit"}\\
7 & Al habilitarse, el texto pasa a blanco bold & Igual que la HU 3.0, criterio 9\\
8 & Al guardar: \code{update} en Room y navegar al Home con la cita modificada & \code{viewModel.\allowbreak{}actualizar(\allowbreak{}cita)} y \code{navigate(\allowbreak{}R.id.action\_editar\_to\_home)}. Si cambiaste la raza, \textbf{regenera la foto} con \code{fotoDeRaza(\allowbreak{}nuevaRaza)} antes del \code{update} (ver el ``Importante'' del enunciado)\\
\end{xltabular}
\end{center}

\begin{notabox}[Lo que el ``Importante'' de la HU 5.0 obliga a hacer]
``Si al momento de editar la cita la raza escogida fue Akita, la imagen debe
mostrar un perro alusivo a esa raza''. Es decir: cambiar la raza \textbf{tiene}
que cambiar la foto del \'item en el Home. Si guardas la misma \code{imagenUrl}
vieja, ese requisito queda incumplido y se nota en la demo (el profesor edita la
raza y compara la foto). Regla: \emph{cada guardado con raza distinta pide una
URL nueva a la API}.
\end{notabox}

% ============================ 8. NAVEGACION ================================
\section{Navegaci\'on: un solo Activity, cinco Fragments}
\label{sec:nav}

El item ``Uso de Fragment'' y el item ``Uso de Navigation'' se califican juntos y
son de los m\'as r\'apidos de verificar: si las pantallas son Activities, o si se
navega con \code{Intent} y \code{startActivity}, se pierden dos items del 60\,\%.

\begin{center}
\small
\begin{xltabular}{\linewidth}{@{}l l l X@{}}
\toprule
\textbf{Desde} & \textbf{Hacia} & \textbf{Acci\'on} & \textbf{Detalle}\\
\midrule
\endhead
\bottomrule
\endlastfoot
\code{loginFragment} & \code{homeFragment} & \code{action\_login\_to\_home} & Solo despu\'es de biometr\'ia exitosa; \code{popUpTo=\allowbreak{}loginFragment}, \code{popUpToInclusive=\allowbreak{}true}\\
\code{homeFragment} & \code{nuevaCitaFragment} & \code{action\_home\_to\_nueva} & FAB rosado\\
\code{homeFragment} & \code{detalleFragment} & \code{action\_home\_to\_detalle} & \code{argument citaId}, tipo \code{integer}\\
\code{nuevaCitaFragment} & \code{homeFragment} & \code{action\_nueva\_to\_home} & \code{popUpTo=\allowbreak{}home} (no deja el formulario en la pila)\\
\code{detalleFragment} & \code{homeFragment} & \code{action\_detalle\_to\_home} & Al eliminar, o con la flecha de atr\'as\\
\code{detalleFragment} & \code{editarFragment} & \code{action\_detalle\_to\_editar} & \code{argument citaId}\\
\code{editarFragment} & \code{homeFragment} & \code{action\_editar\_to\_home} & Tras \code{update} exitoso\\
\code{editarFragment} & \code{detalleFragment} & (bot\'on atr\'as) & \code{navigateUp(\allowbreak{})} al detalle\\
\end{xltabular}
\end{center}

\begin{verbatim}
<fragment android:id="@+id/loginFragment"
    android:name="com.univalle.dogapp.view.login.LoginFragment"
    android:label="Login">
    <action android:id="@+id/action_login_to_home"
        app:destination="@id/homeFragment"
        app:popUpTo="@id/loginFragment" app:popUpToInclusive="true"/>
</fragment>

<fragment android:id="@+id/detalleFragment"
    android:name="com.univalle.dogapp.view.detalle.DetalleFragment">
    <argument android:name="citaId" app:argType="integer"/>
    <action android:id="@+id/action_detalle_to_editar"
        app:destination="@id/editarFragment"/>
</fragment>
\end{verbatim}

\begin{itemize}[leftmargin=1.4em,itemsep=2pt]
  \item \code{MainActivity} contiene \code{NavHostFragment} con
        \code{app:defaultNavHost=\allowbreak{}"true"} y \code{app:navGraph=\allowbreak{}"@navigation/\allowbreak{}nav\_graph"}.
  \item \code{startDestination} = \code{loginFragment}.
  \item Las flechas de las toolbars llaman a \code{navigateUp(\allowbreak{})} o a la acci\'on
        expl\'icita; nunca \code{startActivity}.
  \item Opcional pero recomendable: Safe Args
        (\code{androidx.\allowbreak{}navigation.\allowbreak{}safeargs.\allowbreak{}kotlin}) para no escribir el nombre
        del argumento como texto.
\end{itemize}

% ============================ 9. SPRINT ====================================
\section{Jira y GitHub: el 20\,\% que suele perderse por descuido}

Es frecuente que un equipo tenga la app funcionando y aun as\'i saque 4,0 porque
no puede demostrar gesti\'on. El RA-3 vale tanto como la interfaz completa, y se
construye \emph{durante} el sprint, no la noche anterior.

\subsection{Jira: c\'omo reflejar los criterios}

\begin{center}
\small
\begin{xltabular}{\linewidth}{@{}l X@{}}
\toprule
\textbf{Elemento} & \textbf{Qu\'e crear}\\
\midrule
\endhead
\bottomrule
\endlastfoot
Proyecto & Tipo Scrum, nombre \code{DogApp} o el nombre del equipo\\
\'Epica & \code{DogApp - Sprint 1}. Agrupa las 5 historias\\
Historias & 5 issues tipo \emph{Historia}: \code{HU 1.0 Ventana Login}, \code{HU 2.0 Home}, \code{HU 3.0 Nueva Cita}, \code{HU 4.0 Detalle}, \code{HU 5.0 Editar}. En la descripci\'on, el acta completa y los criterios de aceptaci\'on\\
Criterios & Los criterios del enunciado como \emph{criterios de aceptaci\'on} (checklist) o como subtareas: es lo que el docente revisa\\
Subtareas & Al menos: \code{UI/XML}, \code{ViewModel + LiveData}, \code{Room}, \code{Retrofit}, \code{Navegaci\'on}, \code{QA}. Asignadas a personas distintas\\
Tablero & Columnas \textbf{Por hacer}, \textbf{En curso}, \textbf{Hecho}, con WIP razonable y tarjetas movi\'endose a diario\\
Responsables & Cada historia y subtarea con asignado real (es la evidencia para la nota individual)\\
Evidencia & Capturas del tablero en el PDF de entrega y en el README\\
\end{xltabular}
\end{center}

\begin{notabox}[La prueba que hace el docente]
Abre Jira y mira el \emph{historial de actividad}. Si las 5 historias se movieron
el mismo d\'ia, o si todo lo hizo una sola cuenta, la gesti\'on queda en evidencia
como simulada. Mueve las tarjetas el mismo d\'ia que trabajas y comenta el avance
en la historia (``se implement\'o el autocompletado, falta el filtro por
prefijo''). Eso es exactamente lo que dice la r\'ubrica: ``que le permita construir
aplicaciones de manera colaborativa''.
\end{notabox}

\subsection{GitHub: reglas del repositorio}

\begin{itemize}[leftmargin=1.4em,itemsep=3pt]
  \item \textbf{Nombre:} igual al nombre del equipo (\code{equipoUno},
        \code{equipoDos}, \ldots). El enunciado lo exige dos veces (repositorio y
        archivo PDF).
  \item \textbf{Ramas:} \code{main} (lo entregable) + \code{develop} +
        \code{feature/\allowbreak{}hu1-login}, \code{feature/\allowbreak{}hu2-home},
        \code{feature/\allowbreak{}hu3-nueva}, \code{feature/\allowbreak{}hu4-detalle},
        \code{feature/\allowbreak{}hu5-editar}. Una rama por historia y Pull Request hacia
        \code{develop}.
  \item \textbf{Commits:} convencionales y peque\~nos
        (\code{feat(login):\allowbreak{} biometria con BiometricPrompt},
        \code{fix(home): refrescar lista al eliminar}). Que cada integrante
        commitee con su \textbf{correo de GitHub} configurado: los commits sin
        autor identificable no cuentan como evidencia individual.
  \item \textbf{Pull Requests:} uno por HU, con al menos una revisi\'on de otro
        integrante (\code{Approve}). Esa interacci\'on es la ``evidencia del
        trabajo en equipo''.
  \item \code{.gitignore}: \code{.idea/}, \code{build/}, \code{local.properties},
        \code{*.iml}, \code{.gradle/\allowbreak{}}. \textbf{No} subir \code{local.properties}
        (contiene la ruta del SDK).
  \item \textbf{README.md}: integrantes, qu\'e hace la app, c\'omo compilarla,
        capturas de las 5 pantallas, diagrama de paquetes MVVM, endpoints
        usados, decisiones t\'ecnicas (incluido el tema del endpoint de la foto
        por raza) y enlaces a Jira.
  \item \textbf{Tag} \code{v1.0-sprint1} en el commit de entrega.
  \item \textbf{Congelar el repositorio} antes del 25 de noviembre, 6:30 p.m.:
        cualquier commit posterior cuesta 0,2.
\end{itemize}

% ============================ 10. PLAN =====================================
\section{Plan de ejecuci\'on del sprint}

El sprint va del 11 al 25 de noviembre de 2026, con monitor\'ia el 18. Este orden
es el que menos retrabajo genera: primero el esqueleto que se compila y navega,
despu\'es los datos, despu\'es la red y al final el pulido visual.

\begin{center}
\small
\begin{xltabular}{\linewidth}{@{}l l X l@{}}
\toprule
\textbf{Fase} & \textbf{Fechas} & \textbf{Entregable verificable} & \textbf{Riesgo}\\
\midrule
\endhead
\bottomrule
\endlastfoot
0. Arranque & 11--12 nov & Repo \code{equipoX} creado, Jira con \'epica + 5 historias + subtareas, \code{MainActivity} + \code{NavHostFragment} + 5 Fragments vac\'ios que navegan, temas y colores en \code{values} & Bajo\\
1. Login + marca & 13--14 nov & HU 1.0 completa: fondo, cabeza de perro, fuente, Lottie, biometr\'ia, navegaci\'on con \code{popUpTo}, \'icono de app & Medio (biometr\'ia en emulador)\\
2. Room + Home & 15--17 nov & Entidad, DAO, DB, Repository, ViewModel; Home con RecyclerView leyendo de Room; alta manual de prueba; \textbf{back del Home cierra la app} & Medio\\
3. Formularios & 15--18 nov & HU 3.0 y HU 5.0 con validaciones, m\'aximos, \code{TextInputLayout}, bot\'on que se habilita, mensaje ``Selecciona un s\'intoma'', insertar y actualizar & Alto (estados del bot\'on)\\
4. Monitor\'ia & 18 nov & Demo parcial: login $\to$ home $\to$ nueva $\to$ home. Anotar correcciones del docente & ---\\
5. Retrofit & 19--20 nov & Lista de razas, autocompletado por prefijo con umbral 2, foto por raza con Glide, fallback sin red & Alto (filtro)\\
6. Detalle & 21 nov & HU 4.0 con tarjetas superpuestas, turno, divisor rosado, eliminar y editar & Medio\\
7. Integraci\'on & 22--23 nov & Recorrido completo de las 5 HU, regeneraci\'on de foto al editar raza, tolerancia a fallos de red & Alto\\
8. QA + pulido & 23--24 nov & Checklist de la secci\'on~\ref{sec:checklist} al 100\,\%, README, capturas, PDF de entrega & Medio\\
9. Entrega & 25 nov & PDF \code{equipoX.\allowbreak{}pdf} subido por un solo integrante, tag, repositorio congelado antes de las 6:30 p.m. & ---\\
\end{xltabular}
\end{center}

\begin{techbox}[Reparto sugerido para que todos tengan nota individual]
Con 4 integrantes: \textbf{A} arquitectura + Room + Repository + ViewModel del
Home; \textbf{B} Login + biometr\'ia + sistema de dise\~no (colores, fuentes,
formas, \'icono); \textbf{C} Retrofit + autocompletado + Glide + fallback;
\textbf{D} Detalle + Editar + QA y Jira. Todos commitean, todos mueven tarjetas y
todos pueden explicar la capa que hicieron. Si son 3 o 5, se reparte igual: nadie
queda sin capa propia.
\end{techbox}

% ============================ 11. ERRORES ==================================
\section{Errores que cuestan nota (y su ant\'idoto)}

\begin{center}
\small
\begin{xltabular}{\linewidth}{@{}X X@{}}
\toprule
\textbf{Error} & \textbf{Ant\'idoto}\\
\midrule
\endhead
\bottomrule
\endlastfoot
Cinco Activities en lugar de cinco Fragments: se pierden los items de Fragment y Navigation & \code{MainActivity} \'unica + \code{NavHostFragment} + \code{nav\_graph}. Antes de escribir la primera pantalla, verificar que navega\\
Toda la l\'ogica en el Fragment: se pierde el item MVVM & ViewModel por pantalla; el Fragment solo observa y pinta\\
DAO que devuelve \code{List} en vez de \code{LiveData}: la lista no se refresca al guardar o borrar & \code{LiveData<List<CitaEntity>>} en \code{observarTodas(\allowbreak{})}\\
Llamar a la API en cada tecla del autocompletado & Una carga de razas en \code{init}, filtrado en memoria\\
Mostrar una foto aleatoria aunque se eligi\'o una raza & \code{/breed/\{raza\}/\allowbreak{}images/\allowbreak{}random} y regenerar al editar la raza\\
Guardar el tel\'efono como n\'umero & Guardar como \code{String} con \code{inputType=\allowbreak{}"number"}\\
El bot\'on Guardar arranca habilitado o nunca se habilita & Estado calculado con \code{doAfterTextChanged} sobre los cuatro campos + \code{isEnabled}\\
El bot\'on habilitado no cambia el texto a blanco bold & \code{setTextColor} + \code{Typeface.\allowbreak{}BOLD} al habilitar (criterio 9 de HU 3.0 y criterio 7 de HU 5.0)\\
No hay mensaje ``Selecciona un s\'intoma'' & Validar el \'indice 0 del dropdown antes de guardar\\
El back del Home devuelve al Login & \code{finishAffinity(\allowbreak{})} en el dispatcher + \code{popUpTo} inclusive al entrar al Home\\
Tarjetas del detalle separadas en vez de superpuestas & Margen negativo o anclaje con solapamiento; comprobar en pantalla\\
Toolbars con t\'itulo alineado a la izquierda & T\'itulo centrado seg\'un los pantallazos (Nueva Cita y Detalle)\\
Se olvida el \'icono propio de la app & Adaptive icon \code{ic\_launcher} + \code{android:\allowbreak{}label="dogApp"}\\
La app se cae sin internet & \code{runCatching} + imagen de marcador; el CRUD local debe funcionar offline\\
Falta la animaci\'on Lottie porque el JSON no est\'a en \code{assets} & Copiar el \code{.json} a \code{app/src/main/assets/\allowbreak{}} y usar \code{lottie\_fileName}\\
Commits todos de la misma persona & Reparto por capas + verificar \code{git shortlog -sne} antes de entregar\\
\end{xltabular}
\end{center}

% ============================ 12. CHECKLIST ================================
\section{Checklist de aceptaci\'on}
\label{sec:checklist}

Marcar cada casilla s\'olo despu\'es de verlo en el dispositivo. Es la lista que
se usa como guion de la sustentaci\'on: cada l\'inea es un criterio calificable.

\newpage

\subsection{Arquitectura (RA-1, 60\,\%)}

\begin{center}
\small
\begin{xltabular}{\linewidth}{@{}c X c@{}}
\toprule
& \textbf{Verificaci\'on} & \textbf{Hecho}\\
\midrule
\endhead
\bottomrule
\endlastfoot
1 & Existen los paquetes \code{model}, \code{view}, \code{viewmodel}, \code{repository}, \code{data} & $\square$\\
2 & Ning\'un Fragment importa \code{CitaDao} ni \code{DogApiService} & $\square$\\
3 & Ning\'un ViewModel recibe \code{Context} & $\square$\\
4 & El Repository centraliza Room y Retrofit & $\square$\\
5 & La lista del Home se actualiza sola al insertar, editar y borrar & $\square$\\
6 & Los datos sobreviven al cerrar y reabrir la app (persistencia real) & $\square$\\
7 & Room con \code{@Entity}, \code{@Dao}, \code{@Database} & $\square$\\
8 & Retrofit con \code{@GET}, \code{@Path}, Gson y \code{suspend} & $\square$\\
9 & Las 5 pantallas son Fragments dentro de un \code{nav\_graph} & $\square$\\
10 & \code{citaId} viaja por argumento de navegaci\'on & $\square$\\
11 & Corrutinas en todas las operaciones de datos & $\square$\\
\end{xltabular}
\end{center}

\subsection{Interfaz (RA-2, 20\,\%)}

\begin{center}
\small
\begin{xltabular}{\linewidth}{@{}c X c@{}}
\toprule
& \textbf{Verificaci\'on} & \textbf{Hecho}\\
\midrule
\endhead
\bottomrule
\endlastfoot
12 & Login sin toolbar, fondo \hex{2E2E33}, cabeza de perro con transparencia arriba a la derecha & $\square$\\
13 & ``DogApp'' blanco bold, centrado, con fuente externa & $\square$\\
14 & Huella Lottie animada abajo al centro & $\square$\\
15 & Di\'alogo biom\'etrico con t\'itulo, subt\'itulo y Cancelar & $\square$\\
16 & Home con header oscuro, esquina inferior izquierda curvada, t\'itulo y logo & $\square$\\
17 & \'Items blancos redondeados con elevaci\'on, avatar circular, nombre, s\'intoma y turno \code{\#n} & $\square$\\
18 & FAB rosado de crear abajo a la derecha & $\square$\\
19 & Toolbars oscuras con esquinas inferiores curvadas y flecha rosada & $\square$\\
20 & Campos con hint que flota, borde morado con foco, m\'aximos 15/30/10 & $\square$\\
21 & Lista de s\'intomas con ``S\'intomas'' por defecto & $\square$\\
22 & Bot\'on rosado redondeado con \'icono, inactivo hasta completar, blanco bold al habilitarse & $\square$\\
23 & Detalle con foto superpuesta sobre la tarjeta de datos & $\square$\\
24 & Detalle con turno, raza bold, s\'intoma, l\'inea rosada, propietario y tel\'efono & $\square$\\
25 & FAB de eliminar rosado a la izquierda y de editar azul oscuro a la derecha & $\square$\\
26 & Formulario de edici\'on precargado con los datos del detalle & $\square$\\
27 & \'Icono de la app propio (\code{dogApp}) & $\square$\\
\end{xltabular}
\end{center}
\newpage

\subsection{Funcionamiento y gesti\'on (RA-1 funcional + RA-3)}

\begin{center}
\small
\begin{xltabular}{\linewidth}{@{}c X c@{}}
\toprule
& \textbf{Verificaci\'on} & \textbf{Hecho}\\
\midrule
\endhead
\bottomrule
\endlastfoot
28 & Huella correcta entra al Home; huella fallida muestra error y no entra & $\square$\\
29 & Cancelar sale del di\'alogo y permanece en el login & $\square$\\
30 & El back del Home lleva al escritorio, no al login & $\square$\\
31 & Guardar sin s\'intoma muestra ``Selecciona un s\'intoma'' y no guarda & $\square$\\
32 & Guardar navega al Home y la cita aparece con su turno y su foto de raza & $\square$\\
33 & Editar cambia los datos y la foto si cambia la raza & $\square$\\
34 & Eliminar borra de Room y la cita desaparece del Home & $\square$\\
35 & Con el m\'ovil en modo avi\'on la app no se cae y permite crear citas & $\square$\\
36 & Jira: 5 historias con criterios, subtareas, responsables y movimientos durante el sprint & $\square$\\
37 & GitHub: repo con el nombre del equipo, ramas por HU, PRs revisados, commits de todos & $\square$\\
38 & README con integrantes, capturas, arquitectura y enlaces & $\square$\\
39 & PDF \code{equipoX.\allowbreak{}pdf} con nombres, apellidos, enlace al repo y a Jira & $\square$\\
40 & Repositorio congelado antes de las 6:30 p.m. del 25 de noviembre & $\square$\\
\end{xltabular}
\end{center}

% ============================ 13. ENTREGA ==================================
\section{Entrega y sustentaci\'on}

\begin{center}
\small
\begin{xltabular}{\linewidth}{@{}l X@{}}
\toprule
\textbf{Requisito} & \textbf{Detalle del enunciado}\\
\midrule
\endhead
\bottomrule
\endlastfoot
C\'amara & Todos los integrantes con c\'amara encendida durante la sustentaci\'on. Penalizaci\'on de 0,5\\
Originalidad & Cualquier indicio de copia se penaliza seg\'un los estatutos de la Universidad\\
Formato & Presentaci\'on en equipo, no individual\\
Puntualidad & Hora de inicio: 6:30 p.m.\\
Herramientas & Cada integrante con repositorio, Jira, Android Studio abierto y dispositivo o emulador listo para mostrar la app\\
Subida del PDF & Un solo integrante sube el PDF; el archivo se llama como el equipo (\code{equipoUno.\allowbreak{}pdf})\\
Contenido del PDF & Nombres y apellidos de los integrantes + enlace al repositorio + enlace a Jira\\
Repositorio & Se llama como el equipo (\code{equipoUno})\\
Commits & Ning\'un commit despu\'es de la fecha y hora de entrega. Penalizaci\'on de 0,2\\
\end{xltabular}
\end{center}

\subsection{Guion sugerido de sustentaci\'on (10--12 minutos)}

\begin{enumerate}[leftmargin=1.6em,itemsep=2pt]
  \item \textbf{(30 s)} Abrir Jira y mostrar el tablero con las columnas y las 5
        historias; abrir GitHub y mostrar la pesta\~na de commits y PRs.
  \item \textbf{(2 min)} Login: marca, Lottie, toque, di\'alogo biom\'etrico,
        huella correcta $\to$ Home. Repetir mostrando \emph{Cancelar} y una
        huella fallida.
  \item \textbf{(1 min)} Home: header, \'items con foto de la raza y turno.
        Pulsar atr\'as del tel\'efono y mostrar que sale al escritorio.
  \item \textbf{(2 min)} Nueva Cita: mostrar el bot\'on inactivo, llenar los
        cuatro campos, mostrar el autocompletado con dos letras, dejar el
        s\'intoma vac\'io y provocar ``Selecciona un s\'intoma'', luego guardar y
        ver la cita en el Home.
  \item \textbf{(1,5 min)} Detalle: tarjetas superpuestas y todos los datos.
        Editar: cambiar la raza, guardar y volver a mostrar el Home con la foto
        nueva.
  \item \textbf{(1 min)} Detalle $\to$ eliminar $\to$ Home sin la cita.
  \item \textbf{(2 min)} C\'odigo: abrir el \'arbol de paquetes y mostrar
        Fragment $\to$ ViewModel $\to$ Repository $\to$ DAO/Retrofit. Mostrar
        \code{observarTodas(\allowbreak{})} devolviendo \code{LiveData} y la interfaz de
        Retrofit.
  \item \textbf{(1 min)} Cierre: decisiones t\'ecnicas (endpoint de foto por
        raza, manejo de red ca\'ida, LiveData como fuente \'unica) y reparto de
        trabajo.
\end{enumerate}

\begin{notabox}[Antes de encender la c\^amara]
Compilar en \emph{Release} o al menos en \emph{Debug} limpio, con el emulador ya
encendido y la huella enrolada, la app instalada en el dispositivo y el
repositorio descargado en el port\'atil de cada integrante. La sustentaci\'on
tiene un tiempo fijo: cada minuto gastado en resolver un \code{Gradle sync}
fallido es un criterio que no se alcanza a mostrar.
\end{notabox}

% ============================ ANEXO A ======================================
\newpage
\section{Anexo A: dependencias de \code{build.gradle.\allowbreak{}kts}}

Todas estas l\'ineas se a\~naden al m\'odulo \code{app}, partiendo de la
configuraci\'on de la clase 8 (\code{compileSdk =\allowbreak{} 34}, \code{minSdk =\allowbreak{} 24}, KSP para
Room).

\begin{verbatim}
plugins {
    id("com.android.application")
    id("org.jetbrains.kotlin.android")
    id("com.google.devtools.ksp")
    id("androidx.navigation.safeargs.kotlin")   // opcional pero recomendado
}

android {
    namespace = "com.univalle.dogapp"
    compileSdk = 34
    buildFeatures { viewBinding = true }        // o dataBinding = true
}

dependencies {
    // base
    implementation("androidx.core:core-ktx:1.12.0")
    implementation("androidx.appcompat:appcompat:1.6.1")
    implementation("com.google.android.material:material:1.11.0")
    implementation("androidx.constraintlayout:constraintlayout:2.1.4")

    // MVVM
    implementation("androidx.lifecycle:lifecycle-viewmodel-ktx:2.7.0")
    implementation("androidx.lifecycle:lifecycle-livedata-ktx:2.7.0")
    implementation("androidx.activity:activity-ktx:1.8.2")
    implementation("androidx.fragment:fragment-ktx:1.6.2")

    // Navigation
    implementation("androidx.navigation:navigation-fragment-ktx:2.7.6")
    implementation("androidx.navigation:navigation-ui-ktx:2.7.6")

    // Room
    implementation("androidx.room:room-runtime:2.6.1")
    implementation("androidx.room:room-ktx:2.6.1")
    ksp("androidx.room:room-compiler:2.6.1")

    // Retrofit
    implementation("com.squareup.retrofit2:retrofit:2.9.0")
    implementation("com.squareup.retrofit2:converter-gson:2.9.0")
    implementation("com.squareup.okhttp3:logging-interceptor:4.12.0")

    // Imagenes y animacion
    implementation("com.github.bumptech.glide:glide:4.16.0")
    implementation("com.airbnb.android:lottie:6.3.0")

    // Biometria
    implementation("androidx.biometric:biometric:1.1.0")

    // RecyclerView y CardView
    implementation("androidx.recyclerview:recyclerview:1.3.2")
    implementation("androidx.cardview:cardview:1.0.0")

    // Corrutinas
    implementation("org.jetbrains.kotlinx:kotlinx-coroutines-android:1.7.3")
}
\end{verbatim}

\section{Anexo B: c\'odigo m\'inimo de las tres piezas delicadas}

\subsection*{B.1 Back del Home al escritorio}

\begin{verbatim}
override fun onViewCreated(view: View, savedInstanceState: Bundle?) {
    requireActivity().onBackPressedDispatcher.addCallback(viewLifecycleOwner) {
        requireActivity().finishAffinity()   // cierra la tarea -> escritorio
    }
}
\end{verbatim}

\subsection*{B.2 Boton que se habilita solo y cambia a blanco bold}

\begin{verbatim}
private fun vigilarCampos() {
    val revisar = {
        val listo = etNombre.text.isNullOrBlank().not() &&
                    etRaza.text.isNullOrBlank().not() &&
                    etPropietario.text.isNullOrBlank().not() &&
                    etTelefono.text.isNullOrBlank().not()
        btnGuardar.isEnabled = listo
        btnGuardar.setTextColor(
            if (listo) Color.WHITE else Color.parseColor("#F5A9AD"))
        btnGuardar.setTypeface(null, if (listo) Typeface.BOLD else Typeface.NORMAL)
        btnGuardar.alpha = if (listo) 1f else 0.6f
    }
    listOf(etNombre, etRaza, etPropietario, etTelefono).forEach { campo ->
        campo.doAfterTextChanged { revisar() }
    }
    revisar()
}
\end{verbatim}

\subsection*{B.3 Guardar con foto por raza y navegar}

\begin{verbatim}
btnGuardar.setOnClickListener {
    val sintoma = dropdownSintomas.text.toString()
    if (sintoma == "Sintomas" || sintoma.isBlank()) {
        Snackbar.make(binding.root, "Selecciona un sintoma",
            Snackbar.LENGTH_LONG).show()
        return@setOnClickListener
    }
    viewModel.guardar(
        nombre = etNombre.text.toString().trim(),
        raza   = etRaza.text.toString().trim(),
        dueno  = etPropietario.text.toString().trim(),
        tel    = etTelefono.text.toString().trim(),
        sintoma = sintoma
    )
}

// en el ViewModel
fun guardar(nombre: String, raza: String, dueno: String,
            tel: String, sintoma: String) {
    viewModelScope.launch {
        val url = repository.fotoDeRaza(raza)          // API, con fallback
        repository.guardar(CitaEntity(
            nombreMascota = nombre, raza = raza, propietario = dueno,
            telefono = tel, sintoma = sintoma, imagenUrl = url,
            fechaCreacion = System.currentTimeMillis()
        ))
        _guardado.emit(true)                           // el Fragment navega
    }
}
\end{verbatim}

\begin{notabox}[Nota final de uso de este documento]
Este documento no reemplaza el enunciado: lo descompone. Si hay conflicto entre
un valor de color, un radio o un texto y lo que dice el PDF oficial de las
Historias de Usuario, \textbf{prevalece el PDF oficial}. Los valores de color y
medidas aqu\'i son lecturas razonables de los pantallazos, pensadas para
centralizarse en \code{res/values} y ajustarse una sola vez.
\end{notabox}

\end{document}
